% !TeX root = ../../Book1.tex
% 纲要标题：### 14.3 超越扩张
% 纲要位置：计划纲要.md，第 514 行。

\section{超越扩张}
\label{b1:sec:1403}


% 纲要内容（仅作写作计划，尚非教材正文）：
%
% * 代数独立与超越基
% * 交换引理及超越次数的不变性
% * 有限生成扩张与纯超越扩张
% * 超越次数的塔式公式
%

一个超越元没有单变量多项式关系，几个超越元之间却仍可能有关系。例如 \(t\) 与 \(t^2\) 分别在 \(K\) 上超越，但满足 \(Y-X^2=0\)。要衡量扩张中真正独立的参数数目，必须同时考察有限多个元素的多项式关系。

\ProseParagraphBreak
有限生成扩张可从已给生成元中逐个选取独立参数。一般超越基的存在性则使用 Zorn 引理；比较无限超越基的大小，还需对有限支集作基数计数。

\subsection{代数独立与超越基}
\label{b1:subsec:1403:01}

\begin{definition}\label{b1:def:transcendence-basis}
设 \(L/K\) 为域扩张，\(T\subseteq L\)。如果对任意整数 \(r\geq1\)、任意互异的 \(t_1,\ldots,t_r\in T\) 以及任意 \(f\in K[X_1,\ldots,X_r]\)，等式 \(f(t_1,\ldots,t_r)=0\) 都蕴含 \(f=0\)，则称 \(T\) 在 \(K\) 上\term[daishu-duli]{代数独立}[algebraic independence]。如果 \(T\) 在 \(K\) 上代数独立，并且 \(L/K(T)\) 是代数扩张，则称 \(T\) 为 \(L/K\) 的\term[chaoyue-ji]{超越基}[transcendence basis]。空集视为代数独立，并约定 \(K(\varnothing)=K\)。
\end{definition}

\begin{lemma}\label{b1:lem:independence-adjoining}
设 \(L/K\) 为域扩张，\(T\subseteq L\) 在 \(K\) 上代数独立，且 \(a\in L\setminus T\)。则 \(T\cup\{a\}\) 在 \(K\) 上代数独立，当且仅当 \(a\) 在 \(K(T)\) 上超越。因此，\(L/K\) 的超越基恰是 \(L\) 中按包含关系极大的 \(K\) 上代数独立子集。
\end{lemma}
\begin{proof}
\(K(T)\) 是由 \(T\) 中有限多个元素的多项式比组成的域。若 \(a\) 满足 \(K(T)\) 系数的多项式方程，清除有限多个分母，就得到 \(a\) 与 \(T\) 中有限多个元素之间的非零多项式关系；多项式仍非零，是因为 \(T\) 代数独立。反过来，把一个这样的关系按 \(a\) 的幂收集，至少一个系数是 \(T\) 中元素的非零多项式值，因而非零；这便给出 \(a\) 在 \(K(T)\) 上的代数方程。极大性的等价刻画随即得到。
\end{proof}

对于有限生成扩张，可以直接选择超越基。若 \(L=K(s_1,\ldots,s_n)\)，依次检查这些生成元；只在 \(s_i\) 对已选元素生成的域仍超越时，才把它选入集合 \(T\)。由\cref{b1:lem:independence-adjoining}，\(T\) 始终代数独立。每个未选入的 \(s_i\) 在当时的参数域上代数，因而在最终的 \(K(T)\) 上仍代数；由\cref{b1:lem:finite-algebraic-generators}，它们有限个生成的扩张 \(L/K(T)\) 也代数，所以 \(T\) 是 \(L/K\) 的超越基。一般扩张未必有有限生成集可供逐项检查，下面的存在定理因而采用 Zorn 引理。

\begin{theorem}\label{b1:thm:transcendence-basis}
每个域扩张 \(L/K\) 都有超越基。更一般地，任意在 \(K\) 上代数独立的子集 \(A\subseteq L\) 都可扩充成 \(L/K\) 的超越基。
\end{theorem}
\begin{proof}
考虑包含 \(A\) 的代数独立子集，按包含排序。它包含 \(A\)，因而非空。任一非空链的并仍代数独立：一个多项式关系只涉及有限多个元素，它们已经同属链中的某一个成员。这个并仍包含 \(A\)，故是该链的上界；空链则以 \(A\) 为上界。于是 Zorn 引理（\cref{b1:thm:A-zorn}）给出极大元 \(T\)。由\cref{b1:lem:independence-adjoining}，每个不在 \(T\) 中的元素都在 \(K(T)\) 上代数；\(T\) 中元素本来就在该域中，所以 \(L/K(T)\) 代数。
\end{proof}

\subsection{交换引理及超越次数的不变性}
\label{b1:subsec:1403:02}

\begin{lemma}[交换引理]\label{b1:lem:transcendence-exchange}
设 \(L/F\) 为域扩张，\(x,y\in L\)。若 \(x\) 在 \(F(y)\) 上代数而在 \(F\) 上超越，则 \(y\) 在 \(F(x)\) 上代数。
\end{lemma}
\begin{proof}
把 \(x\) 的一个 \(F(y)\) 系数方程清除分母，得到 \(P(x,y)=0\)，其中 \(P(X,Y)\in F[X,Y]\) 非零。把它按 \(Y\) 展开为 \(\sum_j q_j(X)Y^j\)。由于 \(x\) 在 \(F\) 上超越，每个非零 \(q_j\) 都满足 \(q_j(x)\neq0\)，所以 \(P(x,Y)\) 是 \(F(x)[Y]\) 的非零多项式。它以 \(y\) 为根，故结论成立。
\end{proof}
该结论称为\term[jiaohuan-yinli]{交换引理}[exchange lemma]。它允许用一个新的独立参数替换旧参数，但替换时必须保证新参数在其余已保留参数上超越。

\begin{lemma}\label{b1:lem:finite-transcendence-bound}
设 \(L/K\) 为域扩张，\(b_1,\ldots,b_m\in L\)。若 \(L/K(b_1,\ldots,b_m)\) 代数，则 \(L\) 中任意在 \(K\) 上代数独立的子集至多含 \(m\) 个元素。这里不要求 \(b_1,\ldots,b_m\) 本身代数独立。
\end{lemma}
\begin{proof}
反设 \(a_1,\ldots,a_{m+1}\in L\) 在 \(K\) 上代数独立。设已选取其中的 \(a_1,\ldots,a_r\)，令 \(F=K(a_1,\ldots,a_r)\)，并以 \(D\) 表示尚未替换的 \(b_i\) 构成的集合。归纳假设是 \(\lvert D\rvert\leq m-r\) 且 \(L/E\) 代数，其中 \(E=F(D)\)；\(r=0\) 时这正是题设。令 \(a=a_{r+1}\)。它在 \(F\) 上超越，却在 \(E\) 上代数，因此 \(D\) 非空。在 \(D\) 中取一个包含极小的非空子集 \(C\)，使 \(a\) 在 \(F(C)\) 上代数。任取 \(b\in C\)，若删去 \(b\) 后 \(a\) 仍代数，就与 \(C\) 的极小性矛盾。因此 \(a\) 在 \(F\bigl(C\setminus\{b\}\bigr)\) 上仍超越；这正是交换引理所需的假设。

对底域 \(F\bigl(C\setminus\{b\}\bigr)\) 应用\cref{b1:lem:transcendence-exchange}，得 \(b\) 在 \(F\bigl(C\setminus\{b\},a\bigr)\) 上代数。令
\[
E'=F\bigl(D\setminus\{b\},a\bigr),\qquad
M=F(D,a)=E(a)=E'(b).
\]
由于 \(a\) 在 \(E\) 上代数，\(M/E\) 代数。交换所得的方程在更大的系数域 \(E'\) 上仍有效，故 \(b\) 在 \(E'\) 上代数，即 \(M/E'\) 代数。又因 \(E\subseteq M\subseteq L\)，归纳假设给出 \(L/M\) 代数。旧参数域 \(E\) 与新参数域 \(E'\) 未必互相包含，不能把二者直接排成一条域塔；共同的上层域 \(M\) 给出了如下包含关系：
\[
\begin{tikzcd}[row sep=large, column sep=large]
& L & \\
& M \arrow[u, hook, "\text{代数}"] & \\
E \arrow[ur, hook, "\text{代数}"] && E' \arrow[ul, hook, "\text{代数}"']
\end{tikzcd}
\]
沿右侧的 \(E'\subseteq M\subseteq L\) 应用\cref{b1:thm:algebraic-tower}，得到 \(L/E'\) 代数。这样便用 \(a\) 替换了 \(b\)，同时保持了归纳所需的代数性。每次替换使 \(D\) 少一个元素，故至多 \(m\) 次后它为空。此时 \(L/F\) 代数，但 \(a_{r+1}\) 在 \(F\) 上超越，矛盾。
\end{proof}

\begin{theorem}\label{b1:thm:transcendence-cardinality}
设 \(L/K\) 为域扩张。任意两个 \(L/K\) 的超越基都具有相同基数；这个共同基数称为扩张 \(L/K\) 的\term[chaoyue-cishu]{超越次数}[transcendence degree]，记为 \(\operatorname{trdeg}_K L\)。
\end{theorem}
\begin{proof}
设 \(A,B\) 为超越基。若其中一个有限，不妨设 \(B\) 有限。\cref{b1:lem:finite-transcendence-bound}给出 \(\lvert A\rvert\leq\lvert B\rvert\)，于是 \(A\) 也有限；交换二者再用该引理即得等号。

现在设二者均无限。每个 \(a\in A\) 在 \(K(B)\) 上代数，其一个代数方程只使用 \(B\) 中有限多个元素的有理函数，故可选有限子集 \(B_a\subseteq B\)，使 \(a\) 在 \(K(B_a)\) 上代数。对任意有限集 \(C\subseteq B\)，令
\[
A_C=\{\,a\in A\mid B_a=C\,\},\qquad M_C=K(C)(A_C).
\]
\(A_C\) 的每个元素都在 \(K(C)\) 上代数；由\cref{b1:prop:algebraic-elements-subfield}，\(M_C/K(C)\) 也代数。把\cref{b1:lem:finite-transcendence-bound}用于 \(M_C/K\)，以 \(C\) 为至多 \(\lvert C\rvert\) 个参数，便得到 \(\lvert A_C\rvert\leq\lvert C\rvert\)，因为 \(A_C\subseteq A\) 仍在 \(K\) 上代数独立。

由\cref{b1:prop:finite-support-cardinality}，无限集 \(B\) 的有限子集总数为 \(\lvert B\rvert\)，而每个 \(A_C\) 都有限。因此对应 \(a\mapsto B_a\) 给出 \(\lvert A\rvert\leq\lvert B\rvert\)。交换 \(A,B\) 得反向不等式，再由\cref{b1:prop:cantor-bernstein}得到等号。为各 \(a\) 同时选取有限支集 \(B_a\)，以及这里的无限基数计数，均采用\cref{b1:def:A-choice}中的选择公理；它与选取超越基时使用的选择原则一致。
\end{proof}
\symidx{\(\operatorname{trdeg}_K L\)：扩张的超越次数}
最后一段不能仅用有限次交换代替：无限基的大小是基数问题，有限个参数的局部估计必须结合有限支集的计数，才能得到所需结论。

\subsection{有限生成扩张与纯超越扩张}
\label{b1:subsec:1403:03}

\begin{definition}\label{b1:def:purely-transcendental}
设 \(L/K\) 为域扩张。如果存在 \(K\) 上代数独立的子集 \(T\subseteq L\)，使 \(L=K(T)\)，则称 \(L/K\) 为\term[chun-chaoyue-kuozhang]{纯超越扩张}[purely transcendental extension]。
\end{definition}
对有限集 \(T=\{t_1,\ldots,t_r\}\)，代入给出 \(K(T)\cong K(X_1,\ldots,X_r)\)。一般超越扩张未必纯超越；超越基只保证以参数域 \(K(T)\) 为新基域后，\(L/K(T)\) 是代数扩张。

\begin{proposition}\label{b1:prop:finitely-generated-field-structure}
设 \(L/K\) 为有限生成域扩张。则存在有限超越基 \(T\subseteq L\)，使 \(L/K(T)\) 有限。因而每个有限生成扩张都可以分成一个有限参数的纯超越扩张和其后的有限代数扩张。
\end{proposition}
\begin{proof}
从有限生成集 \(S\) 中按上文的方法逐个选取，得到超越基 \(T\subseteq S\)。每个 \(s\in S\setminus T\) 在 \(K(T)\) 上代数；\(S\setminus T\) 有限，故\cref{b1:lem:finite-algebraic-generators}说明 \(L=K(T)\bigl(S\setminus T\bigr)\) 在 \(K(T)\) 上有限。
\end{proof}

\begin{proposition}[超越次数的塔式公式]\label{b1:prop:transcendence-tower}
设 \(K\subseteq E\subseteq L\) 为域扩张塔。则下式中的加法为基数加法：
\begin{equation}\label{b1:eq:transcendence-tower}
\operatorname{trdeg}_K L=\operatorname{trdeg}_K E+\operatorname{trdeg}_E L.
\end{equation}
\end{proposition}
\begin{proof}
取 \(E/K\) 的超越基 \(T\) 及 \(L/E\) 的超越基 \(U\)。\(T\subseteq E\)，而 \(U\) 的任何元素都不属于 \(E\)，否则它在 \(E\) 上代数；故 \(T\cap U=\varnothing\)。

先证 \(T\cup U\) 在 \(K\) 上代数独立。只涉及 \(T\) 或只涉及 \(U\) 的关系已由各自的独立性排除。一般关系涉及有限多个 \(t_1,\ldots,t_r\in T\) 及 \(u_1,\ldots,u_s\in U\)，按 \(u_j\) 的单项式收集后可写成
\[
\sum_{\beta}P_\beta(t_1,\ldots,t_r)
u_1^{\beta_1}\cdots u_s^{\beta_s}=0,
\qquad P_\beta\in K[X_1,\ldots,X_r],
\]
其中 \(\beta\in\mathbb Z_{\geq0}^s\)，只有有限多个 \(P_\beta\) 非零。各系数 \(P_\beta(t_1,\ldots,t_r)\) 属于 \(E\)；由 \(U\) 在 \(E\) 上代数独立，它们全为零。再由 \(T\) 在 \(K\) 上代数独立，每个 \(P_\beta\) 都是零多项式，故原关系只能是零关系。

又因 \(T\) 是 \(E/K\) 的超越基，\(E/K(T)\) 代数。\(E\) 中每个元素在更大的 \(K(T,U)\) 上仍代数，所以 \(E(U)=K(T,U)(E)\) 在 \(K(T,U)\) 上代数。\(L/E(U)\) 也因 \(U\) 是 \(L/E\) 的超越基而代数；由\cref{b1:thm:algebraic-tower}，\(L/K(T,U)\) 代数。因此 \(T\cup U\) 是 \(L/K\) 的超越基。由于 \(T\cap U=\varnothing\)，其基数为 \(\lvert T\rvert+\lvert U\rvert\)，即得所述等式。
\end{proof}

有限代数扩张的次数沿塔相乘；超越次数沿塔按基数相加。
